\documentclass{amsart}
% For dealing with references we use the comment environment
\usepackage{verbatim}
\newenvironment{reference}{\comment}{\endcomment}
%\newenvironment{reference}{}{}
\newenvironment{slogan}{\comment}{\endcomment}
\newenvironment{history}{\comment}{\endcomment}

% For commutative diagrams we use Xy-pic
\usepackage[all]{xy}

% We use 2cell for 2-commutative diagrams.
\xyoption{2cell}
\UseAllTwocells

% We use multicol for the list of chapters between chapters
\usepackage{multicol}

% This is generally recommended for better output
\usepackage{lmodern}
\usepackage[T1]{fontenc}

% For cross-file-references

% Package for hypertext links:
\usepackage{hyperref}

% For any local file, say "hello.tex" you want to link to please

% Theorem environments.
%
\theoremstyle{plain}
\newtheorem{theorem}[subsection]{Theorem}
\newtheorem{proposition}[subsection]{Proposition}
\newtheorem{lemma}[subsection]{Lemma}

\theoremstyle{definition}
\newtheorem{definition}[subsection]{Definition}
\newtheorem{example}[subsection]{Example}
\newtheorem{exercise}[subsection]{Exercise}
\newtheorem{situation}[subsection]{Situation}

\theoremstyle{remark}
\newtheorem{remark}[subsection]{Remark}
\newtheorem{remarks}[subsection]{Remarks}

\numberwithin{equation}{subsection}

% Macros
%
\def\lim{\mathop{\mathrm{lim}}\nolimits}
\def\colim{\mathop{\mathrm{colim}}\nolimits}
\def\Spec{\mathop{\mathrm{Spec}}}
\def\Hom{\mathop{\mathrm{Hom}}\nolimits}
\def\Ext{\mathop{\mathrm{Ext}}\nolimits}
\def\SheafHom{\mathop{\mathcal{H}\!\mathit{om}}\nolimits}
\def\SheafExt{\mathop{\mathcal{E}\!\mathit{xt}}\nolimits}
\def\Sch{\mathit{Sch}}
\def\Mor{\mathop{\mathrm{Mor}}\nolimits}
\def\Ob{\mathop{\mathrm{Ob}}\nolimits}
\def\Sh{\mathop{\mathit{Sh}}\nolimits}
\def\NL{\mathop{N\!L}\nolimits}
\def\CH{\mathop{\mathrm{CH}}\nolimits}
\def\proetale{{pro\text{-}\acute{e}tale}}
\def\etale{{\acute{e}tale}}
\def\QCoh{\mathit{QCoh}}
\def\Ker{\mathop{\mathrm{Ker}}}
\def\Im{\mathop{\mathrm{Im}}}
\def\Coker{\mathop{\mathrm{Coker}}}
\def\Coim{\mathop{\mathrm{Coim}}}

% Boxtimes
%
\DeclareMathSymbol{\boxtimes}{\mathbin}{AMSa}{"02}

%
% Macros for moduli stacks/spaces
%
\def\QCohstack{\mathcal{QC}\!\mathit{oh}}
\def\Cohstack{\mathcal{C}\!\mathit{oh}}
\def\Spacesstack{\mathcal{S}\!\mathit{paces}}
\def\Quotfunctor{\mathrm{Quot}}
\def\Hilbfunctor{\mathrm{Hilb}}
\def\Curvesstack{\mathcal{C}\!\mathit{urves}}
\def\Polarizedstack{\mathcal{P}\!\mathit{olarized}}
\def\Complexesstack{\mathcal{C}\!\mathit{omplexes}}
% \Pic is the operator that assigns to X its picard group, usage \Pic(X)
% \Picardstack_{X/B} denotes the Picard stack of X over B
% \Picardfunctor_{X/B} denotes the Picard functor of X over B
\def\Pic{\mathop{\mathrm{Pic}}\nolimits}
\def\Picardstack{\mathcal{P}\!\mathit{ic}}
\def\Picardfunctor{\mathrm{Pic}}
\def\Deformationcategory{\mathcal{D}\!\mathit{ef}}

\setlength{\emergencystretch}{3em}
\begin{document}
\title[Stacks Project, Tag 09ZL]{335 --- Finite etale algebras over a henselian pair}
\maketitle
\noindent Copyright (C) 2005--2025 Johan de Jong. Stacks Project authors. Source: \href{https://stacks.math.columbia.edu/tag/09ZL}{Tag 09ZL}.

\noindent Editorial difficulty note (not author text). Difficulty H2: Full faithfulness is proved by identifying algebra maps with sections and then with lifted idempotents. Essential surjectivity uses integral closure, decomposition of the closed fibre, and localization to produce a finite etale lift, so an entire equivalence must be constructed..

\noindent Editorial provenance note (not author text). History: excerpt from revision \texttt{a0444\allowbreak 6e57e\allowbreak c1fbc\allowbreak 252a8\allowbreak 71afc\allowbreak ec775\allowbreak 2fb28\allowbreak 07b14}, more-algebra.tex, lines 3382--3437. Reference presentation was adapted; original-author context and core support blocks were retained or added. The mathematical wording is unchanged. Licensed under GNU FDL 1.2 or later; no invariant sections or cover texts. See ../licenses/COPYING. Context blocks reproduce author definitions and supporting results; they are not additional counted problems.

\begin{definition}
\label{definition-henselian-pair}
A {\it henselian pair} is a pair $(A, I)$ satisfying
\begin{enumerate}
\item $I$ is contained in the Jacobson radical of $A$, and
\item for any monic polynomial $f \in A[T]$ and factorization
$\overline{f} = g_0h_0$ with $g_0, h_0 \in A/I[T]$ monic
generating the unit ideal in $A/I[T]$, there
exists a factorization $f = gh$ in $A[T]$ with $g, h$ monic
and $g_0 = \overline{g}$ and $h_0 = \overline{h}$.
\end{enumerate}
\end{definition}

\begin{lemma}
\label{lemma-finite-etale-equivalence}
Let $(A, I)$ be a henselian pair. The functor $B \to B/IB$ determines
an equivalence between finite \'etale $A$-algebras and finite \'etale
$A/I$-algebras.
\end{lemma}

\begin{proof}
Let $B, B'$ be two $A$-algebras finite \'etale over $A$.
Then $B' \to B'' = B \otimes_A B'$ is finite \'etale as well
(Algebra, Lemmas \href{https://stacks.math.columbia.edu/tag/00U2}{lemma etale (Tag 00U2)} and
\href{https://stacks.math.columbia.edu/tag/02JK}{lemma base change integral (Tag 02JK)}).
Now we have $1$-to-$1$ correspondences between
\begin{enumerate}
\item $A$-algebra maps $B \to B'$,
\item sections of $B' \to B''$, and
\item idempotents $e$ of $B''$ such that $B' \to B'' \to eB''$ is
an isomorphism.
\end{enumerate}
The bijection between (2) and (3) sends $\sigma : B'' \to B'$
to $e$ such that $(1 - e)$ is the idempotent
that generates the kernel of $\sigma$ which exists by
Algebra, Lemmas \href{https://stacks.math.columbia.edu/tag/00U7}{lemma map between etale (Tag 00U7)} and
\href{https://stacks.math.columbia.edu/tag/00U8}{lemma surjective flat finitely presented (Tag 00U8)}.
There is a similar correspondence between
$A/I$-algebra maps $B/IB \to B'/IB'$ and idempotents
$\overline{e}$ of $B''/IB''$ such that
$B'/IB' \to B''/IB'' \to \overline{e}(B''/IB'')$ is
an isomorphism. However every idempotent $\overline{e}$ of $B''/IB''$
lifts uniquely to an idempotent $e$ of $B''$
(Lemma \href{https://stacks.math.columbia.edu/tag/09XI}{lemma characterize henselian pair (Tag 09XI)}).
Moreover, if $B'/IB' \to \overline{e}(B''/IB'')$ is an isomorphism,
then $B' \to eB''$ is an isomorphism too by Nakayama's lemma
(Algebra, Lemma \href{https://stacks.math.columbia.edu/tag/00DV}{lemma NAK (Tag 00DV)}).
In this way we see that the functor is fully faithful.

\medskip\noindent
Essential surjectivity. Let $A/I \to C$ be a finite \'etale map.
By Algebra, Lemma \href{https://stacks.math.columbia.edu/tag/04D1}{lemma lift etale (Tag 04D1)}
there exists an \'etale map $A \to B$ such that $B/IB \cong C$.
Let $B'$ be the integral closure of $A$ in $B$. 
By Lemma \href{https://stacks.math.columbia.edu/tag/09XH}{lemma helper finite type (Tag 09XH)} we have
$B'/IB' = C \times C'$ for some ring $C'$
and $B'_g \cong B_g$ for some $g \in B'$ mapping to $(1, 0) \in C \times C'$.
Since idempotents lift
(Lemma \href{https://stacks.math.columbia.edu/tag/09XI}{lemma characterize henselian pair (Tag 09XI)})
we get $B' = B'_1 \times B'_2$ with $C = B'_1/IB'_1$ and $C' = B'_2/IB'_2$.
The image of $g$ in $B'_1$ is invertible.
Then $B_g = B'_g = B'_1 \times (B_2)_g$ and this implies
that $A \to B'_1$ is \'etale.
We conclude that $B'_1$ is finite \'etale over $A$
(integral \'etale implies finite \'etale by
Algebra, Lemma \href{https://stacks.math.columbia.edu/tag/02JJ}{lemma characterize finite in terms of integral (Tag 02JJ)}
for example)
and the proof is done.
\end{proof}
\end{document}
